\chapter{怎样把素数信息变成零点信息}
\label{chap:bridge}
\section{先把素数计数变得平滑}
定义冯\,芒戈尔特函数
\[
\Lambda(n)=
\begin{cases}
\log p,&n=p^a,\ p\text{ 为素数},\ a\ge1,\\
0,&\text{其他情况}.
\end{cases}
\]
例如 $\Lambda(2)=\log2$、$\Lambda(4)=\log2$、$\Lambda(8)=\log2$，而 $\Lambda(6)=0$。它把素数幂保留下来，正好与 Euler 乘积的对数导数配合。

对实数 $x\ge1$，令
\begin{align}
S(x)&=\sum_{1\le n\le x}\Lambda(n)\left(\frac{x}{n}-1\right),\label{eq:S}\\
G(x)&=S(x)-x\log x+(1+\gamma)x,\label{eq:G}
\end{align}
其中 $\gamma$ 是 Euler 常数，$\log$ 为自然对数。求和范围中的 $n$ 是整数。新项在 $x=n$ 时的权重为零，所以跨过一个素数幂时不会突然给 $S$ 加一个非零跳跃。更精确地，在任意紧的 $x$ 区间内，$S$ 是有限多个连续函数 $\Lambda(n)\max\{x/n-1,0\}$ 的和；因此 $S$、$G$ 及 $g(u)=G(e^u)$ 都连续。

例如
\[
S(4)=\log2+\frac13\log3;
\]
虽然 $4$ 是素数幂，它在此处的权重 $4/4-1$ 恰好为零。

\Expl{平滑不只是为了把图画得好看。权重 $x/n-1$ 的积分可以精确计算，随后会产生一个分母 $s(s-1)$。减去的主项则消掉本来由 $\zeta$ 在 $1$ 的极点造成的干扰。}

\begin{lemma}[无需 RH 的粗增长界]
\label{lem:growth}
存在常数 $C_0>0$，对所有 $x\ge1$ 有 $|G(x)|\le C_0x^2$。
\end{lemma}
\begin{proof}
因为 $0\le\Lambda(n)\le\log n\le n$，求和项非负，且
\[
0\le S(x)\le x\sum_{1\le n\le x}\frac{\Lambda(n)}n
\le x\lfloor x\rfloor\le x^2.
\]
再由 $\log x\le x$ 和 $x\le x^2$，
\[
|G(x)|\le x^2+x^2+(1+|\gamma|)x^2
=(3+|\gamma|)x^2.
\]
可以取 $C_0=3+|\gamma|$。
\end{proof}

\section{亲手算出 Mellin 变换}
先在 $\Rea s>2$ 定义
\[
\Phi(s)=\int_1^\infty G(x)x^{-s-1}\dd x.
\]
上面的粗界保证绝对收敛。对每个正整数 $n$，
\begin{align*}
\int_n^\infty\left(\frac{x}{n}-1\right)x^{-s-1}\dd x
&=\frac1n\frac{n^{1-s}}{s-1}-\frac{n^{-s}}s\\
&=\frac{n^{-s}}{s(s-1)}.
\end{align*}
Euler 乘积在绝对收敛区域给出
\[
-\frac{\zeta'(s)}{\zeta(s)}
=\sum_{n=1}^\infty\frac{\Lambda(n)}{n^s}.
\]
交换求和与积分的绝对收敛性可由
$\sum\Lambda(n)n^{-\Rea s}<\infty$ 及刚算出的实参数积分检验。
同时，
\[
\int_1^\infty x^{-s}\dd x=\frac1{s-1},\qquad
\int_1^\infty(\log x)x^{-s}\dd x=\frac1{(s-1)^2}.
\]
代入定义得到
\begin{equation}
\boxed{\displaystyle
\Phi(s)=-\frac{\zeta'(s)}{s(s-1)\zeta(s)}
-\frac1{(s-1)^2}+\frac{1+\gamma}{s-1}.}
\label{eq:Phi}
\end{equation}
这个等式先在收敛区域成立，右边随后给出 $\Phi$ 的亚纯延拓。

\section{为什么 \texorpdfstring{$s=1$}{s=1} 不留下一个假极点}
$\zeta$ 在 $1$ 的展开为
\[
\zeta(s)=\frac1{s-1}+\gamma+O(s-1).
\]
令 $h(s)=(s-1)\zeta(s)$，在 $1$ 延拓为 $h(1)=1$。则
$h'(1)=\gamma$，且 $h$ 在 $1$ 的一个邻域内没有零点。式\eqref{eq:Phi}可改写为
\begin{equation}
\Phi(s)=-\frac{h'(s)/h(s)+1-(1+\gamma)s}{s(s-1)}.
\label{eq:Phi-regularized}
\end{equation}
分子在 $s=1$ 等于 $\gamma+1-(1+\gamma)=0$，所以分子含有因子 $s-1$。除去它后，$\Phi$ 在 $1$ 是全纯的。

\Expl{如果不先消掉这一项，后面的“发现极点就发现零点”会被 $\zeta$ 本来的极点干扰。式\eqref{eq:G}中的校正项正是为此选择的。}

\section{每个零点都留下一个真正的极点}
设 $\rho\ne0,1$ 是 $\zeta$ 的 $m$ 重零点，$m\ge1$。在它的邻域写
\[
\zeta(s)=(s-\rho)^m a(s),\qquad a(\rho)\ne0.
\]
于是
\[
\frac{\zeta'(s)}{\zeta(s)}=
\frac m{s-\rho}+\frac{a'(s)}{a(s)}.
\]
代入式\eqref{eq:Phi}，得到
\begin{equation}
\Phi(s)=\frac{c_\rho}{s-\rho}+\text{一个在 $\rho$ 附近全纯的函数},
\qquad c_\rho=-\frac{m}{\rho(\rho-1)}\ne0.
\label{eq:residue}
\end{equation}
因此，在开右半平面 $\Rea s>0$，$\Phi$ 的极点恰好就是 $\zeta$ 的零点，且每个都是单极点。零点即使有较高重数，经过对数导数后也成为单极点，重数只进入非零系数 $c_\rho$。

\section{为何可以省掉实极点的特殊分支}
\begin{lemma}
\label{lem:no-real-poles}
$\Phi$ 在正实轴上没有极点。
\end{lemma}
\begin{proof}
当实数 $r>1$ 时，$\zeta(r)=\sum n^{-r}>0$。当 $0<r<1$ 时，考虑配对后的级数
\[
\eta(r)=\sum_{n=1}^\infty\bigl((2n-1)^{-r}-(2n)^{-r}\bigr)>0.
\]
每一项严格为正，且级数收敛。更一般地，在任何紧集 $\Rea s\ge a>0$、$|s|\le M$ 上，
\[
\left|(2n-1)^{-s}-(2n)^{-s}\right|
=\left|\int_{2n-1}^{2n}s t^{-s-1}\dd t\right|
\le M(2n-1)^{-a-1}.
\]
所以配对级数局部一致收敛，定义 $\Rea s>0$ 中的全纯函数。在 $\Rea s>1$ 将偶数项分离，可得
\[
\eta(s)=(1-2^{1-s})\zeta(s).
\]
由恒等定理，该关系延续到 $\Rea s>0$ 且 $s\ne1$。对 $0<r<1$，分母 $1-2^{1-r}<0$，故
\[
\zeta(r)=\frac{\eta(r)}{1-2^{1-r}}<0.
\]
这排除了所有正实零点；$s=1$ 在 $\Phi$ 中又已可去，因此没有正实极点。
\end{proof}

最后令
\begin{equation}
g(u)=G(e^u),\qquad u\ge0.
\label{eq:g}
\end{equation}
换元 $x=e^u$ 后，
\begin{equation}
\Phi(s)=\int_0^\infty g(u)e^{-su}\dd u,
\qquad |g(u)|\le C_0e^{2u},\qquad \Rea s>2.
\label{eq:Laplace}
\end{equation}
这样，零点已被转化为一个有实值输入的拉普拉斯变换的极点。
